\documentclass[11pt]{article}
\usepackage[margin=1in]{geometry}
\usepackage{amsmath,amssymb,amsfonts,bm}
\usepackage{booktabs,multirow,array}
\usepackage{graphicx}
\usepackage{xcolor}
\usepackage{hyperref}
\usepackage{cleveref}
\usepackage{siunitx}
\usepackage{enumitem}
\usepackage{microtype}
\usepackage{authblk}
\emergencystretch=3em
\hypersetup{colorlinks=true,linkcolor=blue!60!black,citecolor=blue!60!black,urlcolor=blue!60!black}
\newcommand{\E}{\mathbf{e}}
\newcommand{\norm}[1]{\left\lVert #1 \right\rVert}
\title{\textbf{CrossPiezo: Databases Agree on Promising Regions but Cannot Resolve the Elite Tail in Piezoelectric Screening}}
\author[1]{Anonymous Author(s)}
\affil[1]{Anonymous Institution}
\date{2 August 2026}
\begin{document}
\maketitle

\begin{abstract}
High-throughput piezoelectric databases are widely used to train and benchmark machine-learning models, yet their cross-source agreement is usually reported as a single global correlation.
Here we reframe cross-database comparison as a question of \emph{screening resolution}: do the Materials Project (MP) and JARVIS databases agree on broad promising regions of chemical space, yet fail to reproducibly resolve the elite tail of candidates?
Using the frozen CrossPiezo-Invariant-v1 benchmark of 573 strict structure-matched pairs (P0) and 207 tighter matches (P2), we compute coordinate-invariant response scalars F1, F3 and F4 and sweep the screened quantile from 1\% to 50\%.
The normalised area under the chance-adjusted concordance curve (nAUCC) for F1 on P0 is 0.105, and the simultaneous 95\% lower confidence bound of adjusted overlap first exceeds 0 at $q=33\%$; no five-quantile persistent onset is observed at $\delta=0.05$, showing that even chance-adjusted agreement is not reliably sustained in the elite tail.
The high-response subset and the tighter P2 panel show even weaker tail resolution, confirming that the disagreement is not a simple sample-size artifact.
Simpler scalar controls (unit-cell volume, band gap) are markedly more cross-source consistent than the piezoelectric response, with positive $\Delta\tau$ and $\Delta$nAUCC.
Source-robust portfolios improve worst-source recall and NDCG relative to JARVIS-only or MP-only selection on a frozen holdout fold, but this is two-source decision robustness, not independent physical validation.
All numbers are hash-bound and traceable to the Phase~7B manifest.
\end{abstract}

\section{Introduction}
Piezoelectric response tensors are central to electromechanical materials design, yet costly to compute with density-functional perturbation theory (DFPT).
The Materials Project (MP) and JARVIS projects have released thousands of computed tensors \cite{dejong2015,choudhary2020}, and recent equivariant neural networks report strong in-source prediction accuracy \cite{yan2024,dong2025,batzner2022}.
Improved agreement with a held-out subset of one database does not guarantee that a candidate is robustly ranked across independently generated high-throughput datasets \cite{hegde2023,wagner2023}.

We introduce \emph{screening resolution} as a diagnostic lens.
Rather than asking ``what is the average rank correlation between two sources?'' we ask ``at what elite fraction $q$ does the overlap between source-specific top-$q$ sets become statistically distinguishable from chance?''
This distinction matters for materials discovery: a model that is validated only on aggregate correlation may still disagree with an independent source on the very candidates a human experimenter would pursue.

\section{Screening-resolution framing}
\label{sec:methods}
\paragraph{Paired universe.}
We use the frozen Phase 6A panel: P0 ($n=573$) and P2 ($n=207$), structure-matched pairs with identical or near-identical crystal structures across MP and JARVIS.
We do not modify panel membership.

\paragraph{Response functionals.}
For each source we compute three coordinate-invariant scalars:
F1 = Cartesian Frobenius norm $\norm{\E}_F$;
F3 = true collinear longitudinal maximum $\max_{\norm{\bm n}=1} |n_i e_{ijk} n_j n_k|$ \cite{maxlongitudinal2021};
F4 = Kelvin/Mandel operator norm on symmetric strain.

\paragraph{Chance-adjusted concordance curve.}
For each panel, metric and quantile $q=1\%,\dots,50\%$, we report the observed top-$q$ Jaccard overlap $J_q$, the exact hypergeometric null expectation $\mathbb{E}[J_q]$ under independent rankings, and the chance-adjusted overlap
\begin{equation}
\widetilde J_q = \frac{J_q - \mathbb{E}[J_q]}{1 - \mathbb{E}[J_q]}.
\end{equation}
We construct a simultaneous 95\% confidence band around $\widetilde J_q$ using a paired bootstrap with a studentized sup-norm critical value \cite{efron1994,romano2005}.
The normalised area under the concordance curve is
\begin{equation}
\mathrm{nAUCC} = \frac{1}{q_{\max}-q_{\min}} \int_{q_{\min}}^{q_{\max}} \widetilde J_q \,dq,
\end{equation}
and the persistent consensus onset is the smallest $q$ for which the lower band exceeds a threshold $\delta$ for at least five consecutive quantiles.

\paragraph{Controls and heterogeneity.}
We compare piezoelectric response against volume, band gap, energy above hull and dielectric trace, recording source fields, common $N$ and available $N$.
We report $\Delta\tau = \tau_{\text{control}} - \tau_{\text{piezo}}$ and $\Delta\mathrm{nAUCC}$ with grouped paired-bootstrap confidence intervals, permutation $p$-values and Benjamini--Hochberg FDR.
Heterogeneity is analysed by crystal system, anion category, response-magnitude quartile, match-quality terciles and heavy-atom fraction, with subgroup minima $N=30$.

\paragraph{Robust portfolio benchmark.}
We fix an elite fraction $q^*=10\%$ and budgets $b \in \{1.0,1.5,2.0\} \times \lfloor q^* n \rfloor$.
Strategies include single-source, average percentile, Borda count, maximin percentile, intersection-first, balanced union, disagreement abstention and a greedy minimax oracle.
We evaluate worst-source recall, worst-source NDCG (with the ideal computed from the full universe), minimax regret and coverage on a deterministic holdout fold obtained by hashing \texttt{pair\_id} modulo 5.

\section{Results}
\label{sec:results}
\subsection{Screening-resolution curves}
\Cref{tab:concordance} summarises the concordance diagnostics.
On P0 the three functionals behave similarly: nAUCC is 0.105 for F1, 0.110 for F3 and 0.106 for F4.
The pointwise lower confidence bound first exceeds 0 at 33\% for F1, but a five-quantile persistent onset at $\delta=0.05$ is not observed for any primary metric, and at $\delta=0.10$ it is not observed for F1.
Thus, even if one screens half of the universe, the two sources only barely agree more than would be expected by chance.

\begin{table}[htbp]
\centering
\caption{Screening-resolution summary for P0 ($n=573$) and P2 ($n=207$). $q_{0.00}$ is the pointwise first crossing; $q_{0.05}^{\mathrm{persist}}$ requires five consecutive lower-confidence bounds above 0.05.}
\label{tab:concordance}
\begin{tabular}{llccc}
\toprule
Panel & Metric & nAUCC & $q_{0.00}$ & $q_{0.05}^{\mathrm{persist}}$ \\
\midrule
P0 & F1\_Frobenius & 0.105 & 33\% & -- \\
P0 & F3\_Longitudinal & 0.110 & 20\% & 38\% \\
P0 & F4\_KelvinOp & 0.106 & 28\% & 42\% \\
P2 & F1\_Frobenius & 0.060 & -- & -- \\
P2 & F3\_Longitudinal & 0.043 & -- & -- \\
P2 & F4\_KelvinOp & 0.046 & -- & -- \\
\bottomrule
\end{tabular}
\end{table}

The high-response subsets (top 50\% by pooled F1) and the tighter P2 panel show \emph{negative} nAUCC values, meaning the elite tail is less reproducible than random ranking.
This confirms that the elite-tail resolution limit is a robust feature of the data, not an artifact of averaging over low-response materials.

\subsection{Scale, order and tail}
Source-wise quantile normalization leaves Kendall $\tau$ essentially unchanged (by monotonic invariance), so the disagreement is not a simple monotonic calibration error.
ECDF Kolmogorov--Smirnov tests and quantile-mapping residuals reveal genuine scale shifts, while top-10\% threshold-crossing counts remain large, indicating that the tail itself is unstable even after an order-only calibration.

\subsection{Property controls}
\Cref{tab:controls} shows that all available controls are more cross-source consistent than piezoelectric F1.
Volume is the strongest control (P0 $\tau=0.967$, nAUCC $=0.918$), followed by band gap ($\tau=0.908$, nAUCC $=0.897$).
All controls have positive $\Delta\tau$ and $\Delta$nAUCC with respect to F1.
The elastic bulk modulus is available only from MP in this processed release, so it is marked \texttt{insufficient\_N} for cross-source comparison.

\begin{table}[htbp]
\centering
\caption{Property-control comparison on P0. Positive $\Delta$ means the control is more consistent than F1.}
\label{tab:controls}
\begin{tabular}{lccccc}
\toprule
Attribute & $\tau$ & nAUCC & $\Delta\tau$ & $\Delta$nAUCC & common $N$ \\
\midrule
F1\_Frobenius & 0.245 & 0.105 & -- & -- & 573 \\
volume & 0.967 & 0.918 & 0.722 & 0.813 & 573 \\
band\_gap & 0.908 & 0.897 & 0.663 & 0.792 & 573 \\
energy\_above\_hull & 0.611 & 0.533 & 0.366 & 0.428 & 573 \\
dielectric\_trace & 0.482 & 0.286 & 0.236 & 0.181 & 573 \\
\bottomrule
\end{tabular}
\end{table}

\subsection{Electronic/ionic decomposition}
Convention-complete electronic/ionic/total tensor decompositions are available for fewer than the required 100 pairs in both sources, so the decomposition analysis is marked \texttt{insufficient\_N}.
We therefore do not draw conclusions about ionic versus electronic contributions.

\subsection{Heterogeneity}
Subgroups with $N\ge 30$ show variation around the panel-average $\tau$, but after FDR correction most subgroup differences are not statistically significant.
Continuous match-quality covariates (RMS distance, lattice distance) are weakly correlated with absolute rank displacement, confirming that the residual disagreement is not dominated by a single match-quality variable.

\subsection{Robust portfolio benchmark}
On the frozen holdout fold, source-robust strategies improve worst-source recall and NDCG relative to JARVIS-only or MP-only selection; the balanced-union strategy reaches a worst-source recall of 1.0 at budget factor 2.0 on P0 F1, while the single-source strategies remain below 0.30 at the same budget.
Disagreement abstention tunes its penalty on the development fold and reduces tail regret by penalising high cross-source disagreement.
Full benchmark numbers are given in the hash-bound \texttt{portfolio\_benchmark.csv} artifact.

\section{Discussion}
The weak elite-tail resolution is a statement about database definitions and processed fields, not about the physical correctness of either source.
Simpler, more convention-robust observables (volume, band gap) are reproduced far more reliably than the piezoelectric tensor, suggesting that a large fraction of the disagreement stems from workflow-specific conventions rather than from irreducible physics.

Source-robust portfolios improve the worst-case recall of a two-source decision maker, but they do not validate physical truth.
A third computational protocol or an experimental adjudication set would be required to identify which source, if either, is ``correct'' on any specific candidate.

\section{Data availability}
Hash-bound result tables, panel membership, the Phase 7B manifest and reproducibility instructions are released in the CrossPiezo repository.

\section{Conclusion}
CrossPiezo shows that MP and JARVIS agree on broad promising regions of piezoelectric chemical space while failing to reproducibly resolve the elite tail.
We recommend that future high-throughput screening reports report screening-resolution curves, chance-adjusted overlaps and robust-portfolio metrics alongside aggregate correlations.

\bibliographystyle{unsrt}
\bibliography{CrossPiezo_ScreeningResolution_references}
\end{document}
